\section{Introduction}

IEEE 802.11ax gives an access point two ways to cut the energy cost of
battery-powered stations. Target Wake Time (TWT) schedules when each station
wakes, so radios sleep between service
periods~\cite{chen2019twt,qiu2021category}. Uplink OFDMA replaces contention
with scheduled resource units, so several stations transmit at
once~\cite{kotagiri2021distributed,han2024multiagent}. Both reduce energy, and
recent work treats their configuration as a scheduling problem of its
own~\cite{busacca2026twt,liu2025priority}.

The two mechanisms are normally configured in isolation. TWT studies fix the
OFDMA setting, and OFDMA studies fix or ignore the sleep
schedule~\cite{karaca2020joint,noh2024joint}. Yet they compete for the same
resource: a trigger frame is only useful to a station awake to receive it, so
every OFDMA round draws on the airtime and the radio states TWT is trying to
conserve.

No single setting serves a whole cell. A basic service set carries mixed
traffic---latency-sensitive video, periodic telemetry and voice
together~\cite{yang2021twt,tan2024drl}---and a schedule tuned for light
telemetry is wasteful under a video burst. A controller must adapt, which
raises how it should be obtained.

A hand-coded rule is cheap to run but needs an expert to tune, and is unlikely
to stay optimal away from the conditions it was tuned for. An idealised online
learner could in principle track the optimum instant by instant, but it must
explore on the live network, and exploration costs latency and risks service
degradation~\cite{wydmanski2021contention,yang2024offline}.

Offline reinforcement learning avoids both
costs~\cite{levine2020offline}. A policy is trained from an already-collected
dataset and never acts on the network while learning, so it needs neither
exploration nor expert-labelled actions---though the state features, reward and
action catalogue remain design choices. The price is
distribution shift, for which conservative value estimation is the standard
remedy~\cite{kumar2020cql,kostrikov2022iql}. Its wireless applications
concentrate on cellular resource
management~\cite{yang2024offline,eldeeb2026marl}, leaving Wi-Fi power save
largely unexamined.

We apply it to joint TWT and OFDMA control. Once per decision window the
access point selects a TWT schedule, a station-to-group assignment and whether
uplink OFDMA is enabled---a 760-way discrete action. Making the OFDMA bit part
of the action rather than a property of the run is what makes the problem
joint. Five offline algorithms are trained on $100{,}235$ transitions logged in
ns-3.44 and compared against hand-designed and non-adaptive baselines.

Our contributions are the following.\footnote{Code, the ns-3 scenario and the
frozen results that generate every table and figure:
\url{https://github.com/shuvo102459/802.11ax-OfflineRL-TWT-OFDMA}}
\begin{itemize}
\item A policy learned without online interaction or expert-set thresholds that
cuts station energy against a tuned hand-designed controller at higher
throughput, holding across ten training runs and against every fixed schedule
tested---its advantage coming from switching schedules, not from finding a
better static one.
\item Evidence that at this operating point enabling the uplink scheduler
raises energy at every TWT configuration tested, with five independently trained
policies ordering by energy as they order by how often each enables it.
\item A service-period-aligned trigger scheduler that solicits a station only
while its service period can contain the whole exchange, removing a race in
which a solicited station sleeps before replying and halving the measured cost
of uplink OFDMA.
\end{itemize}
